\section{Conclusion}
\label{sec:conclusion}

Pangenome graphs have already outgrown the memory of the machines that analyze
them, and compression has already solved their storage problem: a $369$\,GB
interchange file fits in $4.5$\,GB. What has not changed is that every tool
consuming the data still begins by undoing that work, and pays more to rebuild the
expanded form than to compute the answer it was asked for.

We showed that this step is not required by the computation. A compressed
container built against five properties, local expansion, side-table-free symbol
classification, column independence, record independence, and closure under the
domain's symmetries, can serve as the substrate the analyses run on. \sys
evaluates six pangenome analyses as folds over grammar-compressed traversal
columns, including a reference-relative variant projection that we recast as a
topology-only discovery pass followed by a streaming allele fold, and so never
builds the graph that the standard tool holds resident. It is $15$--$613\times$
faster than the established tools with $3$--$25\times$ less memory, and it answers
questions on whole-genome graphs those tools cannot load.

The result we would most like to see generalized is not the speedup but the
inversion behind it. When the compressed form is what the query reads, compression
ratio stops being a tax on access and becomes a multiplier on it, and the incentive
to compress harder aligns with the incentive to compute faster. Getting there costs
a few percent of ratio, spent on keeping the format computable. For any dataset
that will outlive the memory of the machines that read it, that looks like a good
trade.
